\documentclass[UTF8,a4paper,11pt,openany]{ctexbook}
\usepackage{../common/statlearnbook}
\SLSetBookMetadata{第 5 册}{统计学习方法讲义 v2.6.0}{采样方法、主题模型与图排序}
\begin{document}
\setcounter{page}{859}
\setcounter{chapter}{35}
\setcounter{section}{1}
\setcounter{figure}{1}
\pagestyle{headings}
\chaptermark{潜在狄利克雷分配}
\sectionmark{LDA 生成过程与概率图}
板块图中，$N_m$板块嵌在$M$板块内，表示每篇文档的每个词位置都重复生成；
$K$板块表示主题词向量重复$K$次。图\ref{fig:V5-C06-plate-graph}用统一符号展示这一生成依赖。
% v2.6.0 figure UID: FIG-P683-01
% Semantic lock: 18 node / 5 draw, N_m nested in M, K separate; no resizebox.
\tikzset{slfig-FIG-P683-01/.style={font=\fontsize{9.4pt}{11.3pt}\selectfont,every path/.append style={line cap=round,line join=round}}}
\pgfkeysifdefined{/tikz/alt}{}{\tikzset{alt/.code={}}}
\begin{figure}[htbp]
\centering
\begin{tikzpicture}[slfig-FIG-P683-01,slfig high,>=Stealth,
  every node/.style={font=\fontsize{9.4pt}{11.3pt}\selectfont},
  latent/.style={circle,draw=SLTeal,fill=SLTealBG,line width=.80pt,minimum size=8.2mm,inner sep=0pt,font=\fontsize{9.6pt}{11.5pt}\selectfont},
  hyper/.style={draw=none,fill=none,inner sep=1pt,text=SLGray,font=\fontsize{9.6pt}{11.5pt}\selectfont},
  obs/.style={circle,draw=SLBlue,fill=SLBlue,line width=.85pt,text=white,minimum size=8.2mm,inner sep=0pt,font=\bfseries\fontsize{9.6pt}{11.5pt}\selectfont},
  edge/.style={-{Stealth[length=2.1mm,width=1.5mm]},draw=SLBlue,line width=.95pt},
  plate/.style={draw=SLGray!62,line width=.66pt,rounded corners=2.2pt,dash pattern=on 2.6pt off 1.8pt},
  plabel/.style={text=SLGray,font=\fontsize{9.2pt}{11.0pt}\selectfont,inner sep=0pt},
  legendtext/.style={anchor=west,font=\fontsize{9.2pt}{11.0pt}\selectfont,inner sep=0pt},
  alt={对象—关系—结论：完整Bayes LDA盘式图把超参数、潜变量和观测变量分开：每篇文档共享一个主题比例，每个词位拥有一个主题指派，所有文档共享带Dirichlet先验的主题词分布；盘框标明重复次数，箭头只表示条件依赖方向}]
\node[hyper] (alpha) at (-5.95,1.25) {$\alpha$};
\node[latent] (theta) at (-3.05,1.25) {$\theta_m$};
\node[latent] (z) at (-.65,1.25) {$z_{mn}$};
\node[obs] (w) at (1.65,1.25) {$w_{mn}$};
\node[hyper] (beta) at (-5.95,-2.10) {$\beta$};
\node[latent] (phi) at (-.65,-2.10) {$\varphi_k$};

\node[plate,fit=(z)(w),inner xsep=4.6mm,inner ysep=4.4mm] (nplate) {};
\node[plate,fit=(theta)(nplate),inner xsep=5.2mm,inner ysep=6.8mm] (mplate) {};
\node[plate,fit=(phi),inner xsep=4.5mm,inner ysep=3.5mm] (kplate) {};

\draw[edge] (alpha) -- (theta);
\draw[edge] (theta) -- (z);
\draw[edge] (z) -- (w);
\draw[edge] (beta) -- (phi);
\draw[edge] (phi) -- (w);

\node[plabel,anchor=south east] at ($(nplate.north east)+(-1.5mm,1.55mm)$) {$N_m$ 个词位};
\node[plabel,anchor=north east] at ($(mplate.south east)+(0,-1.9mm)$) {$M$ 篇文档};
\node[plabel,anchor=north east] at ($(kplate.south east)+(0,-1.9mm)$) {$K$ 个主题};
\node[obs,minimum size=6.4mm] at (5.45,.95) {};
\node[legendtext] at (6.25,.95) {观测变量};
\node[latent,minimum size=6.4mm] at (5.45,0) {};
\node[legendtext] at (6.25,0) {潜变量};
\node[hyper] at (5.45,-.95) {$\alpha,\beta$};
\node[legendtext] at (6.25,-.95) {超参数（plate 外）};
\end{tikzpicture}
\captionsetup{width=.94\linewidth}
\caption{完整Bayes LDA盘式图把超参数、潜变量和观测变量分开：每篇文档共享一个主题比例，每个词位拥有一个主题指派，所有文档共享带Dirichlet先验的主题词分布；盘框标明重复次数，箭头只表示条件依赖方向}
\label{fig:V5-C06-plate-graph}
\end{figure}

\FloatBarrier
\textbf{学习策略。}已知$\boldsymbol w$后，目标是后验
\[
p(\boldsymbol z,\boldsymbol\theta,\boldsymbol\varphi
\mid\boldsymbol w,\boldsymbol\alpha,\boldsymbol\beta).
\]
对$\boldsymbol z$求和并对两组概率向量积分会产生指数规模耦合，通常不能精确计算。
折叠Gibbs在完整Bayes模型中积分消去共轭概率向量，再对$\boldsymbol z$抽样；
变分EM在点参数变体中，以可计算的平均场分布近似后验并交替更新局部变分参数与全局点参数。
\end{document}
